\chapter{能耗约束下基于可视图的路径规划方法研究}

面向低空物流任务，无人机路径规划是协同规划与控制中的重要环节。规划结果同时受障碍绕行、飞行约束、规划时间和路径质量等因素影响，其目标是在多种约束条件下获得满足规划时间与路径质量要求的路径方案。无人机飞行区域多位于楼宇、厂房和仓储设施密集区，障碍物的关键几何信息主要体现在占地范围和高度。在规划尺度上，这类建筑通常具有较规则的立面、屋顶和边界，因此将城市建筑障碍近似为轴对齐长方体，以统一描述其空间占据范围与高度信息。基于上述环境表示，城市三维路径规划同时面临建筑群造成的绕行约束和建图规模增长问题。可视图能够描述楼宇间的可通行关系，但直接扩展到三维后，节点数量、层间连通关系和候选连边组合迅速增加，建图与搜索开销随之上升。与此同时，仅以欧氏距离作为代价指标时，规划结果难以反映无人机在爬升、平飞、下降和转折阶段的能量差异，而低空物流任务又受电池容量、载荷和航程约束，因此规划目标还需要纳入贴近任务执行过程的飞行动力学代价。

本章建立离散路径点与折线路径模型，提出能耗约束下基于可视图的路径规划方法。该方法在可视图建图阶段结合兴趣区域筛选、分层切片和候选端点筛选，以控制三维建图规模；在路径搜索阶段引入考虑飞行动力学的边权模型，在保证空间可行性的同时统一考虑路径长度与能量消耗。随后在三组不同复杂度场景中，将能耗约束下基于可视图的路径规划方法与基础VG方法及栅格化A*方法进行对比，比较图规模控制、路径质量保持和能量约束表达方面的差异，整体流程如图~\ref{framework1} 所示。

\begin{figure}[!h]
  \centering
  \includegraphics[width=1\textwidth]{pic/算法框图1.png}
  \caption{能耗约束下基于可视图的路径规划方法框架图}
  \label{framework1}
\end{figure}

\section{问题描述}

本章研究静态城市场景中的无人机三维点到点路径规划问题。给定起点 \(s=(x_s,y_s,z_s)\) 和终点 \(g=(x_g,y_g,z_g)\)，需要在由建筑障碍物构成的三维环境中生成一条从起点到终点的可行飞行路径。环境中的建筑障碍物统一采用轴对齐长方体表示，障碍物集合记为 \(\mathcal O=\{O_i\}_{i=1}^{M}\)，其中第 \(i\) 个障碍物可写为 \(O_i=[x_i^{\min},x_i^{\max}]\times[y_i^{\min},y_i^{\max}]\times[z_i^{\min},z_i^{\max}]\)。该表示以边界区间形式刻画障碍物的占地范围与高度，能够较好描述低空物流场景中楼宇、厂房和仓储设施对飞行空域形成的主要几何约束。对于城市低空物流任务而言，路径规划不仅要保证无人机能够绕开障碍物安全抵达目标区域，还需要为后续轨迹生成和控制执行提供具有较好通行性与代价特性的离散航路点序列。

为给飞行过程预留安全裕度，对参与规划的障碍物采用确定性几何膨胀。设膨胀尺度为 \(\eta>0\)，膨胀后的障碍物记为
\begin{equation}
\begin{aligned}
O_i^{+} ={}& [x_i^{\min}-\eta,x_i^{\max}+\eta] \times [y_i^{\min}-\eta,y_i^{\max}+\eta] \\
&\times [z_i^{\min},z_i^{\max}+\eta].
\end{aligned}
\label{eq:2_1}
\end{equation}
障碍物在 \(x\) 和 \(y\) 方向做对称扩张，在 \(z\) 方向仅对顶部进行扩张，以同时考虑机体空间尺寸、路径跟踪误差以及飞行过程中必要的安全间隔。后续可视性判定与路径搜索均基于膨胀后的障碍物集合进行，图~\ref{fig:obstacle-space} 为障碍空间示意。

\begin{figure}[htbp]
    \centering
    \includegraphics[width=0.6\linewidth]{pic/part1/obstacles.png}
    \caption{路径规划障碍空间示意图}
    \label{fig:obstacle-space}
\end{figure}

路径采用由有限个离散路径点连接而成的折线形式，记为 \(\mathcal P=(p_0,p_1,\dots,p_K)\)，其中 \(p_0=s\)、\(p_K=g\)，对应空间折线记为 \(\Gamma(\mathcal P)=\bigcup_{k=0}^{K-1}\overline{p_kp_{k+1}}\)。若对任意 \(k=0,1,\dots,K-1\) 和任意障碍物 \(O_i^{+}\)，均满足 \(\overline{p_kp_{k+1}}\cap O_i^{+}=\varnothing\)，则称路径 \(\mathcal P\) 为可行路径。记全部可行折线路径构成的集合为 \(\Omega\)，则路径规划的目标是在 \(\Omega\) 中寻找总代价最小的解：
\begin{equation}
\mathcal P^{\star} = \arg\min_{\mathcal P\in\Omega} \sum_{k=0}^{K-1} c(p_k,p_{k+1}).
\label{eq:2_2}
\end{equation}
其中 \(c(p_k,p_{k+1})\) 为相邻路径点之间的边代价，用于综合度量飞行距离、姿态变化和能量消耗等因素。进一步地，该问题可写为带权图搜索问题：节点集合由起点、终点以及障碍物边界上提取的候选转折点构成，边集合由空间中彼此可见的节点对构成。由此，连续空间中的路径规划被转化为离散节点序列上的优化问题，输出结果为从起点到终点的路径点序列及其形成的空间折线。

\section{能耗约束下基于可视图的路径规划方法}

三维环境中的基础可视图若直接采用全部障碍物角点并进行全局连边，图规模将迅速膨胀。为控制建图与搜索开销并保留主要绕障拓扑，本章采用兴趣区域筛选、分层切片和极角端点筛选相结合的建图策略，并在图搜索阶段引入综合能耗代价。

\subsection{基于可视图的建图方法}

基础可视图以起点、终点和障碍物顶点作为图节点，任意两节点连线不与障碍物相交时，在二者之间建立可视边。该方法在二维环境中较充分地利用了障碍边界信息，但在三维场景下，障碍物角点数量和节点对组合数同步增加，大量冗余连边将显著抬升可视性检测与建图成本。

建图在 \(xy\) 平面执行兴趣区域筛选。设起终点连线在 \(xy\) 平面中的投影记为 \(\Pi_{xy}(\overline{sg})\)，障碍物 \(O_i\) 的平面投影记为 \(\Pi_{xy}(O_i)\)。保留所有与 \(\Pi_{xy}(\overline{sg})\) 相交的障碍物投影，并对这部分投影顶点求二维凸包，得到与任务航线对应的候选区域 \(\mathcal C\)；再保留全部与 \(\mathcal C\) 相交的障碍物，形成兴趣区域障碍集合。筛选后的建图范围集中于起终点连线附近及主要绕障区域，远离航线的障碍物不再参与节点生成和可视性检测。

兴趣区域确定后，保留障碍物执行固定尺度膨胀，得到用于碰撞检测的障碍物集合，并在此基础上进行分层切片。设膨胀后障碍物上下表面的高度样本集合为 \(\mathcal H_z=\{z_i^{\min},z_i^{\max}\mid O_i^{+}\in\mathcal O^{+}\}\)。从 \(\mathcal H_z\) 中取 \(0.2,0.4,0.6,0.8,1.0\) 五个分位数，并将起点和终点高度 \(z_s,z_g\) 一并纳入，得到水平切层集合 \(\mathcal Z=\mathrm{sort}\{Q_{0.2},Q_{0.4},Q_{0.6},Q_{0.8},Q_{1.0},z_s,z_g\}\)，其中 \(Q_q\) 表示高度样本在分位数 \(q\) 处的取值。切层集合经去重和排序后用于候选节点生成。

对任意 \(z\in\mathcal Z\)，若障碍物 \(O_i^{+}\) 与平面 \(z\) 相交，则在该平面上产生一个矩形截面，其四个角点为
\[
\begin{aligned}
&(x_i^{\min},y_i^{\min},z),\quad (x_i^{\max},y_i^{\min},z),\\
&(x_i^{\max},y_i^{\max},z),\quad (x_i^{\min},y_i^{\max},z).
\end{aligned}
\]
所有切层上的角点与起点、终点共同构成节点集合 \(\mathcal V\)。每个角点同时记录所属障碍物编号，该编号在后续连边阶段用于排除同一障碍物表面的无效贴边连接。

节点生成后，连边阶段不对全部节点对执行全局检测，而在候选邻居筛选后进行可视性判定。对任意当前节点 \(v=(x_v,y_v,z_v)\)，只在与其同层及相邻层的障碍物中寻找候选连边点。若某一障碍物在目标层上的四个角点为 \(u_1,u_2,u_3,u_4\)，则考察由节点 \(v\) 指向四个角点的平面向量在 \(xy\) 平面内的方位顺序，并将四个角点按绕节点 \(v\) 的方位角大小规律排序，仅保留极小角和极大角对应的两个端点作为候选邻居。该筛选保留了绕障过程中更可能形成转折的边界点，同时压缩可视性检测次数。

可视性判定采用基于坐标轴参数区间重叠的线段与长方体相交检测。对任意候选节点对 \((v,u)\)，若两点属于同一障碍物，则不建立连边；否则将线段 \(\overline{vu}\) 表示为 \(p(t)=v+t(u-v)\)，其中 \(t\in[0,1]\)。设膨胀障碍物盒体为 \(B=[x^{\min},x^{\max}]\times[y^{\min},y^{\max}]\times[z^{\min},z^{\max}]\)，并记方向向量 \(u-v=(d_x,d_y,d_z)\)。以 \(x\) 轴为例，若 \(d_x\neq 0\)，则线段落入盒体在 \(x\) 轴投影区间内所对应的参数范围为
\[
\left[\min\left(\frac{x^{\min}-x_v}{d_x},\frac{x^{\max}-x_v}{d_x}\right),\quad
\max\left(\frac{x^{\min}-x_v}{d_x},\frac{x^{\max}-x_v}{d_x}\right)\right].
\]
若 \(d_x=0\) 且 \(x_v\notin[x^{\min},x^{\max}]\)，则线段不可能与盒体相交。对 \(y\) 轴和 \(z\) 轴作同样处理。仅当三条坐标轴对应的参数区间与 \(t\in[0,1]\) 存在公共重叠部分时，线段才进入盒体内部，因而判定为与障碍物相交；若任一坐标轴上的参数区间不重叠，则线段不与该盒体相交。对全部膨胀障碍物均满足不相交时，在节点 \(v\) 与 \(u\) 之间建立无向可视边。为减小数值误差带来的边界误判，计算过程中对盒体边界引入极小收缩量。

起点与终点进一步进行一次直连检测。若线段 \(\overline{sg}\) 与任一膨胀障碍物均不相交，则直接在起点与终点之间建立连边。由上述步骤得到的图结构写为 \(G=(\mathcal V,\mathcal E)\)，其中 \(\mathcal V\) 为节点集合，\(\mathcal E\) 为可视边集合。上述建图过程可整理为算法~\ref{alg:improved_evg}。

\begin{algorithm}[H]
\caption{基于可视图的建图流程}
\label{alg:improved_evg}
\DontPrintSemicolon
\SetAlgoLined
\LinesNumbered
\KwIn{起点 \(s\)，终点 \(g\)，障碍物集合 \(\mathcal O\)，膨胀尺度 \(\eta\)}
\KwOut{带权图 \(G=(\mathcal V,\mathcal E, W)\)}
\BlankLine
根据连线投影 \(\Pi_{xy}(\overline{sg})\) 筛选相交障碍物，由其投影顶点构造凸包 \(\mathcal C\)\;
对 \(\mathcal O\) 中与 \(\mathcal C\) 相交的障碍物按尺度 \(\eta\) 执行三维膨胀，生成集合 \(\mathcal O^{+}\)\;
由 \(\mathcal O^{+}\) 的上下表面高度及 \(z_s, z_g\) 共同构造切层高度集合 \(\mathcal Z\)\;
提取 \(\mathcal O^{+}\) 在所有 \(z \in \mathcal Z\) 截面上的矩形四角点，与 \(s, g\) 构成节点集 \(\mathcal V\)\;
初始化边集 \(\mathcal E \leftarrow \emptyset\)\;
\ForEach{节点 \(v \in \mathcal V\)}{
    在同层与相邻层中提取各障碍物极小与极大角端点，构成候选邻居集 \(\mathcal N_v\)\;
    \ForEach{候选节点 \(u \in \mathcal N_v\)}{
        \If{参数区间检测显示线段 \(\overline{vu}\) 与 \(\mathcal O^{+}\) 无干涉}{
            \(\mathcal E \leftarrow \mathcal E \cup \{(v, u)\}\)\;
        }
    }
}
\If{线段 \(\overline{sg}\) 与 \(\mathcal O^{+}\) 无干涉}{
    \(\mathcal E \leftarrow \mathcal E \cup \{(s, g)\}\)\;
}
利用代价函数计算 \(\mathcal E\) 中各边权重，返回带权图 \(G=(\mathcal V, \mathcal E, W)\)\;
\end{algorithm}

\subsection{能耗代价与路径搜索}

旋翼机飞行能耗建模通常从推进功率出发，将旋翼诱导功率、叶型功率、机体寄生阻力功率以及高度变化引起的势能项纳入统一表达。Zeng等针对旋翼无人机通信轨迹优化问题建立了包含叶型功率、诱导功率和寄生功率的推进能耗模型\cite{zeng2019energy}；Kirschstein面向无人机包裹配送任务进一步分析了飞行模式、环境条件和悬停过程对能耗的影响\cite{kirschstein2020comparison}；Yan等则给出了适用于旋翼无人机推进能耗估计的简化闭式模型\cite{yan2021new}。基于上述研究，本文在图边代价计算中采用由诱导功率、机体阻力功率和爬升功率组成的飞行能耗表达。

设图中一条边连接节点 \(p_i\) 和 \(p_j\)，边长记为 \(d_{ij}=\|p_j-p_i\|_2\)。在固定巡航速度 \(v\) 下，该边对应的飞行时间为 \(t_{ij}=d_{ij}/v\)，两端节点高度差对应的爬升角为 \(\theta_{ij}=\arcsin\left(\frac{z_j-z_i}{d_{ij}}\right)\)。设无人机总质量为 \(m\)，重力为 \(W=mg\)，空气密度为 \(\rho\)，等效气动阻力面积为 \(C_DA\)，则机体阻力可写为 \(D=\frac{1}{2}\rho C_DA v^2\)，维持该飞行状态所需的推力近似为
\[
T=\sqrt{W^2+D^2+\rho C_DA v^2W\sin\theta_{ij}}.
\]
若总旋翼盘面积记为 \(A_r\)，则诱导速度为 \(w=\sqrt{T/(2\rho A_r)}\)。在此基础上，边 \((i,j)\) 上的总功率表示为
\begin{equation}
P_{ij} = \kappa wT + \frac{1}{2}\rho C_DA v^3 + Wv\sin\theta_{ij}.
\label{eq:2_3}
\end{equation}
其中 \(\kappa\) 为诱导功率修正系数。于是该边的飞行能量为 \(E_{ij}=P_{ij}t_{ij}\)。

路径搜索采用距离项与能量项加权的综合代价：
\begin{equation}
c(i,j) = k_1E_{ij} + k_2d_{ij},
\label{eq:2_4}
\end{equation}
其中 \(k_1\) 和 \(k_2\) 分别表示能量项和距离项的权重。路径总代价沿各边逐段累加，得到折线路径 \(\mathcal P\) 的总成本。

路径搜索在带权图 \(G=(\mathcal V,\mathcal E)\) 上采用A*算法。对任意节点 \(n\)，累计已走代价记为 \(g(n)\)，启发项取为节点 \(n\) 到终点 \(g\) 的直连综合代价估计 \(h(n)=c(n,g)\)，搜索优先级函数为 \(f(n)=g(n)+h(n)\)。
A*算法在优先队列上维护待扩展节点，每次取 \(f(n)\) 最小的节点进行扩展，并在发现更低累计代价时更新其前驱节点。当终点首次从优先队列中弹出时，依据前驱信息回溯得到从起点到终点的离散路径点序列，并输出对应的空间折线路径。

\section{对比实验}

为验证所提方法在城市低空物流典型点到点运输任务中的适用性，构造三组三维静态障碍场景开展数值仿真对比实验。仿真任务可理解为无人机从一处配送节点飞往另一处任务节点，飞行过程中需要在城市建筑群之间完成安全绕障与能耗受限飞行。为保证不同方法比较时任务条件一致，三组场景的起点和终点均固定为 \(s=(0,0,10)\) 和 \(g=(780,780,40)\)。地图边长约为 \(920\,\mathrm{m}\)，障碍物均采用轴对齐长方体表示，其在 \(xy\) 平面内的尺寸范围为 \(30\sim100\,\mathrm{m}\)，高度范围为 \(10\sim50\,\mathrm{m}\)，用于模拟厂房、仓储设施和城市楼宇对低空通行形成的三维约束。

考虑到城市区域内建筑分布密度和通行通道复杂度会直接影响可视图建图规模与搜索效率，实验分别设置 20、40 和 80 个障碍物，构造低、中、高三种复杂度场景，以考察方法在不同环境复杂度下的适应性。在上述统一场景设置下，对比对象包括能耗约束下基于可视图的路径规划方法、基础VG方法和栅格化A*方法。基础VG方法直接使用障碍物角点构造三维可视图并进行全局连边，栅格化A*方法采用 \(5\,\mathrm{m}\) 分辨率和 26 邻域搜索。三种方法在相同场景下均采用一致的综合代价模型进行评价，对比指标为运行时间、路径长度和能量消耗。

与能耗计算有关的主要参数设定如表~\ref{tab:drone-params} 所示。参数取值与仿真中的边权模型保持一致。

\begin{table}[htbp]
    \centering
    \caption{仿真参数设定}
    \label{tab:drone-params}
    \begin{tabular}{p{0.20\textwidth}p{0.12\textwidth}p{0.16\textwidth}p{0.18\textwidth}p{0.20\textwidth}}
        \hline
        参数名称 & 符号 & 数值 & 单位/说明 & 备注 \\
        \hline
        机体质量 & \(m\) & 50 & kg & 默认总质量 \\
        巡航速度 & \(v\) & 12 & m/s & 边代价计算速度 \\
        空气密度 & \(\rho\) & 1.225 & kg/m$^3$ & 标准大气条件 \\
        阻力系数 & \(C_D\) & 0.3 & 无量纲 & 等效阻力系数 \\
        机体参考尺寸 & \(l\times w\) & 0.6 \(\times\) 0.6 & m \(\times\) m & 机身平面尺寸 \\
        旋翼半径 & \(R\) & 0.3 & m & 单旋翼半径 \\
        旋翼数量 & \(N\) & 4 & 个 & 四旋翼配置 \\
        诱导功率系数 & \(\kappa\) & 1.15 & 无量纲 & 功率修正系数 \\
        障碍膨胀尺度 & \(\eta\) & 2 & m & 几何安全裕度 \\
        \hline
    \end{tabular}
\end{table}

\begin{figure}[htbp]
    \centering
    \includegraphics[width=1\linewidth]{pic/part1/visible.png}
    \caption{40 个障碍物场景下三种方法规划结果可视化对比}
    \label{fig:visible-compare}
\end{figure}

\noindent 图~\ref{fig:visible-compare} 为 40 个障碍物场景下三种方法的规划结果可视化。能耗约束下基于可视图的路径规划方法与基础VG方法所得路径均沿障碍物边界组织绕行，栅格化A*方法的局部转折更密集，路径形态呈现更明显的阶梯特征。



图~\ref{fig:20compare}、图~\ref{fig:40compare} 和图~\ref{fig:80compare} 分别为三组场景下三种方法的结果。三幅图对应不同的随机障碍场景，不同场景之间的绝对数值不要求随障碍物数量严格单调变化，结果分析以同一场景内各方法的相对表现为主。

\begin{figure}[htbp]
    \centering
    \includegraphics[width=0.5\linewidth]{pic/part1/20compare.png}
    \caption{20 个障碍物场景下三种方法的性能对比}
    \label{fig:20compare}
\end{figure}



在 20 个障碍物场景中，能耗约束下基于可视图的路径规划方法的运行时间最短，但路径长度和能量消耗略高于基础VG方法。该场景障碍物分布较稀疏，主要通行通道数量有限，基础VG方法形成的稠密可视图尚未带来明显计算负担，同时保留了更多跨障碍物连通关系和绕行选择，所得路径在几何上更直接。能耗约束下基于可视图的路径规划方法采用候选节点压缩和候选边筛选后，计算量显著下降，但稀疏场景中仍可能进入最优解的部分连边被同步删除，路径质量因而出现小幅损失。该场景下，图规模压缩的收益主要体现为时间开销下降，路径质量优势尚未形成。

\begin{figure}[htbp]
    \centering
    \includegraphics[width=0.5\linewidth]{pic/part1/40compare.png}
    \caption{40 个障碍物场景下三种方法的性能对比}
    \label{fig:40compare}
\end{figure}


在 40 个障碍物场景中，能耗约束下基于可视图的路径规划方法在运行时间、路径长度和能量消耗三项指标上均优于其余两种方法。障碍数量增加后，三维可视图中的冗余连边显著增多，基础VG方法虽然保持了充分连通性，但大量边与最终较优路径无关，建图和搜索过程逐步受到无效空间关系主导。能耗约束下基于可视图的路径规划方法将建图范围限制在任务航线附近和主要绕障区域，通过分层切片维持必要的垂向连通，再利用极角端点筛选保留有效转折点，使图结构更接近实际绕障拓扑。图规模压缩同时降低了计算开销与无效扩展数量，路径质量和运行效率同步改善。

\begin{figure}[htbp]
    \centering
    \includegraphics[width=0.5\linewidth]{pic/part1/80compare.png}
    \caption{80 个障碍物场景下三种方法的性能对比}
    \label{fig:80compare}
\end{figure}

在 80 个障碍物场景中，能耗约束下基于可视图的路径规划方法的优势进一步扩大。高复杂度环境下，基础VG方法的节点对数量和可视性检测次数快速增长，边集规模急剧膨胀，计算资源大量消耗在冗余连边的生成与判定上。过密图结构进一步引入更多低价值候选边，削弱了搜索对关键绕障通道的聚焦能力，最终表现为更高的时间开销与总代价。能耗约束下基于可视图的路径规划方法在兴趣区域筛选阶段剔除了远离主航路的障碍物，在垂向上仅保留必要切层，在各障碍物局部仅保留最可能形成转折的端点，因此在复杂环境中保持了持续的图规模控制能力和较稳定的路径质量。

综合三组实验结果，能耗约束下基于可视图的路径规划方法的运行时间始终最短，所采用的图压缩策略表现出稳定的规模控制能力。路径质量随场景复杂度呈现明显分化：在 20 个障碍物的稀疏场景中，基础VG方法凭借更充分的全局连通关系在路径长度和能量消耗上略占优势；当障碍物数量增加至 40 和 80 时，能耗约束下基于可视图的路径规划方法在两项指标上均优于基础VG方法和栅格化A*方法。该方法在中高复杂度环境中更有效地平衡了图规模、搜索效率与路径质量。

与栅格化A*方法相比，能耗约束下基于可视图的路径规划方法在三组场景中的路径长度和能量消耗均更低。栅格化A*方法的搜索性能直接受分辨率约束：分辨率较细时，状态空间显著增大，搜索时间迅速上升；分辨率较粗时，连续空间中的几何细节被离散近似，路径转折受限于有限邻接方向，容易形成明显的阶梯状折线。该离散误差同时增加路径长度和姿态调整次数，并在综合代价中累积为额外能量消耗。能耗约束下基于可视图的路径规划方法直接利用障碍物边界构造候选连边，在连续几何结构上组织搜索，所得路径转折更少、通行方向更协调。

\section{本章小结}

本章围绕低空物流场景中的无人机三维路径规划问题，建立了以离散路径点和折线路径为基础的规划模型，并据此构建用于路径搜索的三维可视图。兴趣区域筛选压缩了参与建图的障碍物集合，分层切片和极角端点筛选进一步降低了冗余节点与连边数量，基于飞行动力学的边权模型将路径长度与能量消耗统一纳入搜索过程。对比实验表明，能耗约束下基于可视图的路径规划方法在城市三维障碍环境中较好地兼顾了建图规模、搜索效率与路径质量，在连续空间表达、冗余连边抑制和能量约束建模方面保持了稳定的综合优势，并输出可作为低空物流协同规划输入的离散路径点，为后续协同轨迹规划与控制提供路径参考。
